\section{关系 III}

\subsection{序关系}

\subsubsection{偏序关系}

\begin{definition}[偏序关系]
    集合 $A$ 上的二元关系 $R$ 称为 $A$ 上的\textbf{偏序关系}，当且仅当 $R$ 是自反的、反对称的、传递的。
\end{definition}

对于一个 $A$ 上的偏序关系 $R$，可以用 $\langle A, \preceq \rangle$ 来表示，称为一个\textbf{偏序结构}。$\langle x, y \rangle \in R$ 可以记做 $x \preceq y$。若对于 $x, y \in A$ 有 $x \preceq y$ 或 $y \preceq x$，那么称 $x, y$ 在偏序关系 $R$ 中是\textbf{可比的}。

比如，自然数域和小于等于关系 $\langle \NN, \le \rangle$ 是偏序结构，集合和包含于关系 $\langle \mathcal{P}(A), \subseteq \rangle$ 是偏序结构。

\subsubsection{全序关系}

\begin{definition}[全序关系]
    设 $\langle A, \preceq \rangle$ 是一个偏序结构，若 $\forall\,x, y \in A: x \preceq y \lor y \preceq x$，那么称 $\preceq$ 为 $A$ 上的\textbf{全序}或\textbf{线序}，称 $\langle A, \preceq \rangle$ 为\textbf{全序结构}或\textbf{链}。
\end{definition}

比如，自然数域和小于等于关系 $\langle \NN, \le \rangle$ 是全序结构，但集合和包含于关系 $\langle \mathcal{P}(A), \subseteq \rangle$ 不是全序结构。

\subsubsection{严格偏序关系}

\begin{definition}[严格偏序关系]
    $R$ 是集合 $A$ 上的\textbf{严格偏序关系}当且仅当 $R$ 是反自反的和传递的。
\end{definition}

比如，自然数域上的小于关系是严格偏序关系。
